\section{Conclusion}
Exact eligibility pooling converts heterogeneous service allocation into a correction that an ordered command path pays once. It gives an exact finite-catalog joint algorithm when eligibility is common, even with a variable symbol budget. The selected-prefix method complements that result: it handles arbitrary catalog eligibility in polynomial time at every fixed symbol budget, and its price and free-completion bounds certify the original problem at every interruption. The two regimes should be selected according to the structure actually present, not presented as a universal computational ranking.

Catalog refinement can increase a sufficient search dimension without changing the winning policy. Our refinement theorem, threshold-cell analysis, affected-history guarantee, and certified elimination rule make this dependence explicit. The experiments show both nontrivial strict-prefix scaling and exact closure of the inherited small hard instances, while retaining all interrupted runs and their measured gaps. These are reproducible computational observations, not a hardness classification or an application calibration.

An installed alphabet measures the commands available at the terminal interface. It does not count target tables, encoder logic, read-only constants, probability precision, or certificate storage. Those resources remain separate, and verification costs are reported rather than treated as free. The preserved mathematical archive contains the broader institutional, continuous-design, robustness, and resource-augmentation results under their original assumptions; the current article develops a coherent exact joint-design argument without deleting that record.
